\section{Introduction}
Evaluating a language model on a benchmark with thousands of items is expensive when repeated for many checkpoints, prompts or decoding settings. A tempting remedy is to run the new (``target'') model on a small subset and extrapolate to the full benchmark, using the item-level records of earlier (``source'') models to choose the subset and to calibrate the extrapolation. Many methods follow this recipe---clustering-based anchor points, IRT-based imputation, learned assessors, active item selection~\citep{vivek2024anchor,polo2024tiny,kipnis2025metabench,pacchiardi2024100,li2025active}---and report accurate predictions from a small fraction of the items. A recent large comparison cautions that these gains are largest when the target resembles the source models, and shrink or vanish otherwise~\citep{zhang2025misses}. Whether to trust a predicted score therefore depends on how the target relates to the source population, which is exactly what is unknown for a new model.

We study this question with a simple frame: the benchmark is a finite population of items, the target's accuracy is a population mean, the coreset is a sample, and source models supply auxiliary information. In this frame a method has a \emph{design} (which items) and an \emph{estimator} (how to turn the target's coreset responses into a score), and the two can be varied independently; the classical survey-sampling toolkit---stratification, Horvitz--Thompson weights, Neyman allocation, model-assisted residual correction---gives a vocabulary and a set of strong ``simple'' baselines that predictor-based work often omits.

\paragraph{Contributions.}
(i) A controlled comparison on 13 binary item-level benchmarks (7 Open LLM Leaderboard, 6 HELM) of designs and estimators, under three kinds of source/target relationship (random split, held-out organisations, stronger targets) at five budgets, with paired hierarchical-bootstrap uncertainty. (ii) Two small ingredients: \emph{VA-PPS}, a variance-aware unequal-probability design for residual-corrected estimators, and \emph{Cal-IRT} / \emph{ShrinkCal-IRT}, per-target logit recalibration of an IRT predictor (with a pretest shrinkage between uncalibrated and recalibrated predictions). (iii) A negative result for tuning by pseudo-shift validation within the source models. (iv) Complete code to regenerate every table and figure.
